%% ---------------------------------------------------------------------
\section{Dos sistemas con varios equilibrios}
\label{sec:bi-ejemplos}
%% ---------------------------------------------------------------------

Los dos ejemplos de esta sección tienen dos equilibrios estables separados por una silla. Esa configuración, la \clave{biestabilidad}, es la forma más simple de memoria que tiene un sistema físico: el estado final depende de dónde empezó.

\subsection{La ecuación de Duffing}

\begin{historia}[Georg Duffing (1861--1944)]
Duffing fue un ingeniero mecánico alemán, formado en Karlsruhe, que se dedicó por su cuenta al estudio de las vibraciones. En 1910 hizo una pasantía en la Westinghouse, en Estados Unidos, y en 1913 se estableció en Berlín como inventor independiente. En 1918 publicó \emph{Erzwungene Schwingungen bei veränderlicher Eigenfrequenz und ihre technische Bedeutung}~\citep{Duf18}, sobre oscilaciones forzadas con frecuencia natural variable, donde aparece la ecuación que lleva su nombre. Con dificultades económicas después de la guerra, en 1921 entró a dirigir el laboratorio de la Ölwerke Stern-Sonneborn (Ossag) en Hamburgo, donde trabajó en la viscosidad y la fricción de los lubricantes~\citep{KoBr11}. Su ecuación se volvió famosa medio siglo después, cuando la versión forzada se convirtió en uno de los ejemplos clásicos de caos.
\end{historia}

La forma general de la ecuación es
\[
  \ddot x+\delta\dot x+\alpha x+\beta x^3=\gamma\cos(\omega t),
\]
donde $\delta\ge0$ es el amortiguamiento, $\alpha$ la rigidez lineal, $\beta$ la no linealidad del resorte ($\beta>0$, resorte que se endurece; $\beta<0$, que se ablanda) y $\gamma\cos\omega t$ el forzamiento. Sin forzamiento y con $\alpha=-1$, $\beta=1$, describe una viga de acero entre dos imanes: un pozo de potencial doble. Con $x_1=x$, $x_2=\dot x$,
\begin{equation}
  \dot x_1=x_2,\qquad \dot x_2=x_1-x_1^3-\delta x_2 .
  \label{eq:bi-duffing}
\end{equation}

\begin{ejemplo}[label=ej:bi-duffing]{Duffing amortiguado}
Halle y clasifique los equilibrios de \eqref{eq:bi-duffing} para $\delta>0$.
\solucion $x_2=0$ y $x_1(1-x_1^2)=0$: equilibrios en $(0,0)$ y $(\pm1,0)$. La jacobiana es $\bigl[\begin{smallmatrix}0&1\\1-3x_1^2&-\delta\end{smallmatrix}\bigr]$.
\begin{itemize}
  \item $(0,0)$: $\delta_A=\det A=-1<0$, silla, con $\lambda=\tfrac12\bigl(-\delta\pm\sqrt{\delta^2+4}\bigr)$.
  \item $(\pm1,0)$: $\tau=-\delta<0$, $\det A=2>0$, $\tau^2-4\det A=\delta^2-8$. Son \emph{focos} estables si $\delta<2\sqrt2\approx2.83$ y \emph{nodos} estables si $\delta\ge2\sqrt2$.
\end{itemize}
Con $\delta=0.5$: silla con $\lambda=0.781$ y $-1.281$, focos con $\lambda=-0.25\pm1.392j$. Los tres son hiperbólicos.
\end{ejemplo}

La energía $E=\tfrac12x_2^2-\tfrac12x_1^2+\tfrac14x_1^4$ (cinética más potencial de doble pozo) cumple $\dot E=-\delta x_2^2\le0$: el amortiguamiento disipa energía, y casi todas las trayectorias terminan en el fondo de uno de los dos pozos. ¿En cuál? La respuesta la da la silla y se ve en la \cref{sec:bi-variedades}: desde $(1.5,0)$ la viga termina en $(1,0)$, pero desde $(2,0)$, que tiene más energía, cruza la barrera y termina en $(-1,0)$.

\codigorepo{cap04_bidimensionales/C4_Ej9_Duffing.m}{duffing}{Equilibrios de Duffing en función de $\delta$ y dos condiciones iniciales vecinas}

\begin{trampa}[label=tr:bi-caos]{``La ecuación de Duffing sin amortiguamiento ni forzamiento es caótica''}
\plausible{Duffing es el ejemplo clásico de caos, y sin fricción nada frena las oscilaciones.}
Con $\delta=\gamma=0$ el sistema es plano e invariante en el tiempo, y conserva $E$: cada trayectoria vive en una curva de nivel de $E$, y el retrato es un ocho rodeado de órbitas cerradas (\cref{ej:oc-duffing}). Las trayectorias de un sistema autónomo en el plano no pueden ser caóticas: el teorema de Poincaré--Bendixson (\cref{cap:orbitas}) solo deja equilibrios, órbitas cerradas y ciclos de equilibrios como comportamiento a largo plazo.
\correcto{El caos de Duffing aparece con el forzamiento $\gamma\cos\omega t$: el sistema forzado tiene un estado de dimensión tres ($x_1$, $x_2$ y la fase del forzamiento), y en tres dimensiones Poincaré--Bendixson ya no aplica.}
\end{trampa}

\subsection{El diodo túnel con su inductancia}

En la \cref{sec:un-diodo} se estudió el circuito del diodo túnel sin inductancia. Con la inductancia $L$ en serie con la fuente, el estado son el voltaje del condensador y la corriente del inductor (\cref{prob:in-diodo}), y el modelo de \citet[\S1.2.2]{Kha02} es
\begin{equation}
  \dot x_1=\frac1C\bigl(-h(x_1)+x_2\bigr),\qquad \dot x_2=\frac1L\bigl(-x_1-Rx_2+E\bigr),
  \label{eq:bi-diodo}
\end{equation}
con $x_1=v_C$ en V, $x_2=i_L$ en mA, $h$ la característica de la \cref{sec:un-diodo}, $E=1.2$\,V, $R=1.5$\,k$\Omega$, $C=2$\,pF, $L=5$\,$\mu$H y $t$ en ns.

Los equilibrios cumplen $x_2=h(x_1)$ y $x_2=(E-x_1)/R$: son los mismos tres cortes con la recta de carga del capítulo~2, $Q_1=(0.063,\,0.758)$, $Q_2=(0.285,\,0.610)$ y $Q_3=(0.884,\,0.210)$. La jacobiana es
\[
  A=\begin{bmatrix}-h'(x_1)/C&1/C\\-1/L&-R/L\end{bmatrix},
  \qquad
  \det A=\frac{1+Rh'(x_1)}{CL},
\]
de modo que la silla aparece exactamente donde $h'(x_1)<-1/R$, la misma condición del capítulo~2. Numéricamente:
\begin{center}\small
\begin{tabular}{lccc}
\toprule
 & $h'(\eq_1)$ & $\lambda_1,\lambda_2$ & tipo\\
\midrule
$Q_1$ & $7.24$ & $-3.59,\ -0.33$ & nodo estable\\
$Q_2$ & $-3.64$ & $1.77,\ -0.25$ & silla\\
$Q_3$ & $2.87$ & $-1.34,\ -0.40$ & nodo estable\\
\bottomrule
\end{tabular}
\end{center}
El circuito es biestable: es una celda de memoria. La inductancia no cambia los equilibrios ni su estabilidad, pero sí el \emph{camino}: en el modelo de primer orden el estado se mueve sobre la recta de $v_C$; aquí cada trayectoria es una curva en el plano, y cuál de los dos nodos se alcanza ya no depende solo de $v_C(0)$, sino también de $i_L(0)$.

\begin{figure}[ht]
\centering
\includegraphics[width=0.85\linewidth]{cap04/f_diodo}
\caption{Retrato del diodo túnel \eqref{eq:bi-diodo}. Las nulclinas son la característica $x_2=h(x_1)$ (coral) y la recta de carga (ámbar). En negro, la variedad estable de la silla $Q_2$: separa las condiciones iniciales que terminan en $Q_1$ de las que terminan en $Q_3$.}
\label{fig:bi-diodo}
\end{figure}

\codigorepo{cap04_bidimensionales/C4_Ej10_DiodoTunel.m}{diodo}{Equilibrios, jacobianas y retrato del diodo túnel de segundo orden}

\begin{pausa}
En el diodo túnel con $v_C(0)=0.4$\,V, a la derecha de $Q_2$, ¿termina siempre el circuito en $Q_3$, como en el modelo de primer orden?
\respuesta{No. Con $i_L(0)=0$ termina en $Q_1$: la corriente del inductor arranca en cero, el diodo descarga el condensador y $v_C$ cruza hacia la izquierda de $Q_2$ antes de que $i_L$ alcance a sostenerlo. Hace falta $i_L(0)\gtrsim0.189$\,mA para terminar en $Q_3$. Al revés, con $v_C(0)=0.15$\,V (a la izquierda de $Q_2$) e $i_L(0)=1.5$\,mA se termina en $Q_3$. La frontera entre las dos cuencas es la variedad estable de $Q_2$ (\cref{fig:bi-diodo}), y depende de los dos estados.}
\end{pausa}
